\documentclass{article}
\usepackage{amsmath}
\usepackage{hyperref}
\usepackage[a4paper]{geometry}
\usepackage{fancyhdr}
\pagestyle{fancy}
\lhead{Mutter-Tocher-Zerfall}
\rhead{Februar 2026}
\begin{document}
\section{Mutter-Tocher-Zerfall}
Zerfällt ein radioaktiver Mutterstoff in einen Tochterstoff, welcher auch radioaktiv ist, so ist dies der \emph{Mutter-Tochter-Zerfall}. Der Tocherstoff kann in ein stabilen Stoff zerfallen.
 
Die Aktivität des Tocherstoffs kann dabei schneller ansteigen als die Aktivität des Mutterstoffs sinkt. 
 
\subsection{Als Tabellenkalkulation}
Dem Prinzip Tabellenkalkulation der \hyperref[Aktivität]{Aktivität} nach, nur mit mehr Spalten, kann auch der Verlauf der Stoffmengen, der Aktivitäten der Stoffe und der Gesamtaktivität des Mutter-Tocher-Zerfalls berechnet werden.
 
Die $N$ und $A$ Spalten werden jeweils in zwei geteilt, eine für die Mutter, eine für die Tochter. Zudem ist ab nun $N_t(i) = N_t(i-1) - A_t(i-1) + A_m(i-1)$
 
\subsection{Verlauf mit $\lambda$} 
Der Verlauf der Aktivität eines Stoffes hängt natürlich an erster Stelle von dessen $\lambda$, aber auch von dessen Menge, ab. Je größer ein $\lambda$, desto höher ist dessen Aktivität, sinkt aber auch wieder schneller. Je größer die Aktivität des Mutterstoffes, desto schneller steigt vorerst die Menge des Tocherstoffes.
\end{document}